% Includable body. Uses only standard LaTeX commands and math.
% Editorial synthesis of the supplied 19 October 2026 volume; no new native run.
\section*{Closing synthesis: established results and the physical joins ahead}
\addcontentsline{toc}{section}{Closing synthesis: established results and the physical joins ahead}

The strength of this volume lies in keeping a source, its receiving law, and the result of that law in the same account. Its first chapters contain substantial mathematical and computational developments: a complete interacting source--field action, finite-window clock certificates, source-sensitive retention and nonclosure results, spatial receiving constructions, controlled signal and record protocols, and strong-field conditional theorems. Those results survive even where the physical identification of an apparatus, geometry, or cosmic inventory remains unfinished. The later chapters extend the account to thermal boundaries, cosmology, halos, radiation, and the Home continuation. Their standing depends on the additional premises at each extension.

This synthesis replaces the earlier closing review and its interface appendix as the current reading guide. Their chapter-item identifiers are retained below so that no question disappears merely because the edition changes. References identify chapters and sections rather than assembled pages, which changed when the expanded Chapters I--VI were incorporated. Historical numerical outputs are reported from their source records; the present editorial synthesis neither reruns nor newly certifies them. Repository-only claims remain source-reported until their precise hypotheses, operands, and evidence are incorporated.

\subsection*{1. A common standard of claims}

A result is \emph{derived within a model} when it follows from its stated premises and domain. A measured value is \emph{recovered under an adapter} when the receiving law supplies the structure responsible for agreement. An expression is \emph{selected or matched} when comparison helped choose it. A \emph{prospective prediction} requires an observable, selection rule, uncertainty treatment, and failure criterion fixed before its test data are examined. A \emph{physical closure} additionally supplies the interface that makes the model quantity the measured one. These descriptions can coexist: an exact conditional calculation can be mathematically complete while its physical coupling remains open.

A certificate concerns its declared preparation, observable, interval, domain, and error allowance. A finite-time estimate is not automatically an asymptotic decay law. A positive global energy form is not a pointwise positive stress tensor. A coordinate singularity is not an observer-independent obstruction to crossing. A retained classical exterior is not a dynamical theory for rotating, moving, or self-gravitating sources. A local encoded record is not an intrinsic remote-field measurement. A finite thermal kernel is not an absolute luminosity without a populated source and a realization clock.

The distinction between numerical error and model uncertainty is equally important. Outward arithmetic enclosures can be very narrow without determining physical calibration, experimental systematics, or the uncertainty associated with a selected formula. Likewise, propagating the measured CMB temperature through a fixed cosmological cascade produces a conditional input-error budget, not a model-independent uncertainty on the Hubble constant. Discrete choices require explicit accounting; they need not be misrepresented as ordinary continuous fit parameters.

\subsection*{2. What the source chapters establish and retain}

\paragraph{Chapter I: the actual interacting clock.}
The signed Write construction, all six driven lanes, the complete source energy, two real field polarizations, both counterterms, and the correlated hard preparation must remain together. Chapter I, Sections 6--7, establishes positive actual phase rates on the declared one-bare-cycle interval and transports the complete action into that phase coordinate. The Berry one-form is retained. The rate is history-dependent, and its inverse is a pointwise Jacobian, not an independently chosen control or the inverse accumulated phase. Neither six lane certificates nor their verification counts establish arbitrarily long autonomous operation. The source-to-occupied-path and physical-unit joins remain separate questions.

\paragraph{Chapter II: a receiver-dependent retained account.}
Present output, predictive state, timed history, and energetic complement are different objects. Equal cardinality between two ledgers does not identify their states or supply a transition map. The finite source and paid reference protocols establish operational distinctions under declared receivers. Nonclosure ranks and finite invariant upper constructions must retain their actual scope; a computed lower bound is not a proven minimal invariant closure. The source-to-space extension requires a complete energy receiver and a physical incidence, including the commuting even energy, interaction, field, reference, and counterterms.

\paragraph{Chapter III: spatial response with its coupling exposed.}
The Poisson equation and test-body law explicitly import the Newtonian coupling and the passive-to-inertial mass assignment. Their Green functions, source profiles, finite boundaries, energy accounts, and continuum errors are substantive consequences of those assignments. Static potential and an endpoint clock do not determine the relativistic spatial response. The microscopic localization, physical ruler, moving-source completion, and spatial metric coefficient cannot be supplied merely by renaming a retained count as a density.

\paragraph{Chapter IV: complete histories and actual readout.}
The uniform complete lapse transports the full Hamiltonian through one internal-age map. Local control and donor preparations, unequal accepted branches, paired polarization records, complete failure accounting, and switched-reference constructions are developed further in Sections 79--82. These current results take precedence over the explicitly retained predecessor dependency ledger. Section 82 supplies exact uniform switched-reference and spatial-residual machinery but does not thereby evaluate the physical nonuniform cap mean, angular noise, or external vacuum-work account. Huge finite preparation budgets remain existence constructions under specified apparatus laws. Appendix I distinguishes archived external-contact receipts from new experimental reductions; its binary-pulsar receipt is not an independent timing fit.

\paragraph{Chapter V: strong-field statements with explicit domains.}
The source lift reaching unity is not a statement that the floor plus lift reaches unity. Static adapter results, approach theorems, curved-port protocols, recovered thermal kernels, and finite area-entry algebra have distinct premises. The source-law discussion, including the radial null-energy-condition alternative, must govern every later reference to infinite approach or sealing. Section 13 and the subsequent developments retain actual source, preparation, backaction, field complement, and resource interfaces. A comparator that supplies an observable orbit does not become a selected SAM stress law.

\paragraph{Chapter VI: records with their physical payments.}
The record and accessibility results use ordinary quantum states and allowed channels. The shared-field moving-source and coherent detector constructions retain all outcomes and physical complements. The final ``Completed and remaining physical objects'' statement distinguishes a paid, co-located encoder from a primitive route-selection mechanism and from an intrinsic remote port. Autonomous centre-of-mass dynamics, microscopic routing, calibration, and resource realization remain to be supplied. Later references to information custody must respect this channel dependence rather than imply a new probability law or unobserved signalling.

\subsection*{3. The cross-volume physical interfaces}

Five interfaces recur. First, one physical unit chart must connect the model's time, ruler, action, and energy to measured quantities without changing calibration between tests. Second, a self-consistent curved source--field action and renormalized stress must determine geometry and motion, rather than receive them separately. Third, a boundary condition for the actual two-polarization field must connect thermal and ringdown problems. Fourth, microscopic control and apparatus laws must realize the already specified donor, routing, pointer, vacuum-ramp, and reset accounts. Fifth, long-time preparation, replenishment, and cosmological evolution must connect finite-source calculations to macroscopic rates.

These are research interfaces, not permission to erase completed conditional work. A physical unit chart would close a dimensional join but would not derive a spatial metric. A boundary condition would define a spectral problem but would not alone specify its excitation or luminosity. A radius mechanism from a different source lineage can be investigated as a candidate coupling; identifying its radius with Chapter I's coupling radius requires a map, not numerical proximity.

For gravity, endpoint redshift, light propagation, static second-order motion, frame dragging, orbital radiation, preferred-frame effects, and clock-species universality are separate exams. A Schwarzschild exterior recovers its test-body observables by assumption. It does not establish every parametrized post-Newtonian parameter of a full many-body theory, nor strong-equivalence-principle behavior of self-gravitating bodies. A binary radiation-flux calculation must carry actual source motion, tensor coupling, conservation, and physical normalization into its comparison.

For cosmology, an inventory identity, a background evolution equation, a recombination calculation, and a likelihood comparison are also separate. Numerical proximity of selected parameters to posterior centroids neither derives their physical allocation nor makes the full spectrum test vacuous. It does mean that spectrum agreement cannot be counted as independent evidence for every formula chosen using those same data. Independent instruments help test transportability and systematics, but do not erase selection history.

\subsection*{4. The chapter register}

The identifiers below correspond to the earlier review's register. A retained question is described at its present scope; an acknowledged answer is not listed again as though absent. Additional interfaces from the earlier Chapters I--VII appendix are carried inside the relevant paragraphs.

\subsubsection*{Chapter I: source preparation, isolation, and calibration}

\paragraph{I.1--I.2: isolation and radiation.}
The review attributes a hard-branch isolation proof to the resolved-clock record H002471, cited by Chapter I. Incorporation requires its exact charge-sector decomposition, thresholds, multiplicities, and separation inequalities. A complement form bound and trial subspace must be used with the relevant min--max statement; a summary alone does not authenticate a new continuum certificate. Existing Chapter I provenance remains intact. The review also proposes a golden-rule coefficient from Chapter IV, Section 77. Its quoted $3.7\times10^{-41}$ per model time is a reported asymptotic estimate, not a newly established relaxation law. Uniform control of the large-time limit, depletion, repeated preparation, and conversion to physical time are still needed. A finite emitted-number bound cannot close those questions or exclude every cosmological source term.

\paragraph{I.3--I.6: physical source joins.}
The retained source radius near $0.1635$ and the separate radius-mechanism value near $0.151$ must be identified by lineage and meaning before comparison. The complete-Hamiltonian coupling principle remains distinct from changing only one source coefficient. One physical calibration must fix time and energy, with a ruler where spatial observations are used. The occupied-path relation between the original phase ledger and quantum generator must carry common phase, commuting even operator, and preparation. The original connector-to-occupied-source-to-joint-field incidence, physical selection of $s$ and $\alpha$, modified vacuum, hard-Hamiltonian control port, and actual timing reference remain explicit interfaces. Exact inverse and error-budget constructions do not constitute achieved hardware calibration.

\subsubsection*{Chapter II: the floor and retained state}

\paragraph{II.1--II.2: Berry content and floor normalization.}
The joint Berry contribution must be evaluated on the clock-active preparation before it can be replaced by a reduced phase formula. The floor drops out of static gradients and shared linear differences, but that does not make it absent from every nonlinear receiver. For the specifically normalized raw-floor square root, a fixed independent gravitational calibration would produce the review's extra linear factor $1/(1-A_0)$. Redshift experiments can test that specified hypothesis. They do not uniquely derive subtraction before every square root, prove universality across all clocks, or prove that the floor is observable only cosmologically. Any claimed exclusion must state the mass, force, and potential calibration held fixed.

\paragraph{II.3--II.4 and retained interfaces.}
The relation among cadence, interval, distance, and a physical metric remains to be written with its time and length units. Preferred-frame tests require a specified preferred-frame dynamics, not merely the existence of a substrate vocabulary. A physically smeared energy density needs declared incidence for source, field, interaction, boundary, support, and reference. The paid Write/inverse/wait protocol still needs a physical realization. The actual frustrated-spin invariant closure retains its nonclosure bounds and finite upper construction without an unproved minimum. These questions were explicit in the earlier interface appendix and are preserved here.

\subsubsection*{Chapter III: static gravity and its dynamical extension}

\paragraph{III.1--III.4.}
The spatial-response coefficient $\kappa_s$ and its connection to $\gamma$ remain a decisive physical join. The review's criterion $\gamma=m-1$ follows only after its chosen radius--lapse, carrier-speed, and metric identifications. Applying it to another source family requires common definitions, units, route cells, and regime of validity. It is a conditional diagnostic, not a completed derivation or falsification of all SAM carriers. The static Schwarzschild adapter supplies the corresponding test-body second-order lapse and perihelion behavior; a full Nordtvedt prediction needs more than that exterior. A velocity-coupled/vector sector remains necessary for dragging. The nonuniform source energy/stress form, original-state remainder, apparatus errors, spatial metric response, moving-source completion, and common SI chart retain their Chapter III and later-chapter scopes.

\subsubsection*{Chapter IV: signals, switching, and finite resources}

\paragraph{IV.1--IV.3.}
The sign of an isolated susceptibility must not be inferred from a total response dominated by uniform time change. The raw-floor comparison has the conditional status stated under II.2. Section 77's finite radiation result remains valuable at its established interval; the golden-rule suggestion does not replace the open long-time physical relaxation and cosmological amplitude rows. The positive Euclidean rate measure and its transform also need a detector and a mapping to cosmological variables before becoming a distance-rate law.

\paragraph{IV.4--IV.6.}
The distance-ratio material should be identified by its receiving assumptions and original provenance; editorial relocation must preserve its actual mathematics and archived benchmark distinctions. Preparation counts of order $10^{58}$ are finite existence budgets, not feasible apparatus proposals. The present Sections 79--82 must be used to update the quantitative nonuniform cap problem: actual unequal branches, incoming hard field, transported shell return, both signed insertions, full angular noise, and paired records belong in the same estimate. Uniform switched existence and a local record do not replace the missing nonuniform output bound. Microscopic donor/controller realization, common units, throughput and reset, vacuum-ramp work, and long-time repeated preparation remain separate physical obligations.

\subsubsection*{Chapter V: approach, boundary, and source law}

\paragraph{V.1--V.3.}
The source-coordinate edge, the finite-impulse conditional theorem, and comparator orbit calculations are already present. Their questions should not be reopened as though the chapter omitted them. The retained resolution is conditional: a source satisfying the theorem's radial energy-condition hypotheses is constrained by it; otherwise the actual violation or flux must be calculated and propagated. No universal infinite-cost crossing rule follows independently. Comparator photon circles and innermost stable circular orbits are observables of their specified metric, not derived physical stress assignments.

\paragraph{V.4--V.5.}
A rotating extension is required before importing Kerr comparisons into Chapter XI. The source-to-curved-port coupling must retain support, angular complement, initial dressing, counterterms, reservoir, switching, backaction, and reference energy. A curved source action with renormalized original-state stress must select the traveler and boundary supply. Autonomous recharge, repetition, selection, and stopping remain open. The reported roughly fifteen-percent amplitude--distance interval can be sharpened within its source certificate; repetition alone does not remove that systematic interval.

\subsubsection*{Chapter VI: accessibility and physical records}

\paragraph{VI.1--VI.3.}
A second-order shared-field velocity/dephasing coefficient must be computed from the actual propagators if it is to replace an existing bound; a bound is not its coefficient. The chapter already states its use of ordinary quantum probability. The remaining physical questions are microscopic routing and control, source-profile calibration using the derived stability radius, autonomous centre-of-mass or boost-deformed action, unforced arm geometry and source-clock preparation, a genuinely remote intrinsic source observable with propagation and extraction, and all vacuum-ramp, actuator, bath, and reset payments. Conditional erasure and feed-forward are different experiments and cannot silently alter the all-attempt account.

\subsubsection*{Chapter VII: thermal scale, entropy, and information}

\paragraph{VII.1--VII.2.}
The temperature factor ledger already identifies the static metric, normalized Killing time, Hawking conversion, and retained quarter-area premise. Preserve it prominently. Recovering the Bekenstein--Hawking first law under these premises is distinct from deriving the quarter-area allocation from the microscopic source. Exact finite-entry and Catalan-series consequences remain established conditional mathematics. Their last-entry regime should be compared with discrete-area proposals without claiming an absolute emission rate.

\paragraph{VII.3--VII.5.}
Relative entropy is stationary at a smooth reference state, so comparing its first derivative with a nonstationary area observable does not test every proposed area--entropy relation. The informative Si29 comparison must align perturbation order and compare Hessians under the same variables and normalization. Nonzero Jacobian enclosures establish local invertibility at their reference points; an explicit quantitative neighborhood and physical boundary-area map remain to be supplied. A material surface and a causal horizon have different emission and accretion implications; surface-flux constraints require a specified model before application. The two-polarization boundary condition, source-selected boundary state, transition strengths, realization clock, depletion, energy join $\epsilon_Q=\hbar\omega_{\rm phys}$, species sum, and outward/inward correlations are still the route from finite kernel to radiation law.

\subsubsection*{Chapter VIII: inventory and Hubble cascade}

\paragraph{VIII.1--VIII.2.}
The finite grammar and its selection history must accompany each cosmological identity. The earlier review's broad enumeration is a sensitivity exercise until its script, candidate rules, duplicates, target windows, and normalization are auditable. Candidate counts alone are not calibrated probabilities of accidental agreement. The baryon allocation factor $(1-\chi)$, exponent $\mathcal L/V=6$, and weight $7/4$ require physical allocation or preparation laws; reusing upstream integers does not alone select their cosmological meaning.

\paragraph{VIII.3--VIII.4.}
The clean $\Omega_m=1/\pi$ and effective inventory must each have an operator, background role, and comparison protocol fixed before use. One cannot switch between them after seeing which agrees better. A frozen Hubble ledger should list thermal anchor, constants, density convention, radiation and neutrino treatment, and all external inputs. The quoted central value near $67.2504\,\mathrm{km\,s^{-1}\,Mpc^{-1}}$ is conditional on that cascade. The historical temperature-only error near $0.021$ is not total model uncertainty or a claim of superior cosmological precision. Literature comparisons marked ``verify'' in the earlier review are not retained as established exclusions.

\paragraph{VIII.5--VIII.6.}
A nucleosynthesis calculation from the selected baryon-to-photon ratio can yield useful independent light-element tests, with nuclear, thermal, and observational errors kept separate. A standard network would test the packet under standard thermal physics, not constitute native substrate nucleosynthesis. The complement $1-\Omega_m$ does not by itself determine a stress tensor or an equation of state. A cosmological-constant adapter and a dynamical $w(z)$ model must be distinguished. Accumulation-based distance or redshift mechanisms must make explicit predictions for spectral redshift, time dilation, temperature evolution, surface brightness, and distance relations before being compared with observations.

\subsubsection*{Chapter IX: recombination and spectrum comparison}

\paragraph{IX.1--IX.3.}
The native recombination lane and the standard-CAMB full-shape lane answer different questions and must remain separate. The packet was frozen before spectrum generation, but its selectors had already been compared with Planck parameters. Report that chronology. The full-shape outputs are source-reported comparisons, not new reductions in this edition. The earlier review identifies a likely second-peak comparator error: the draft's $520$ must not be retained without checking the cited peak definition and dataset. Use one peak finder and windows on both spectra. Report sample sizes, windows, covariance, nuisance handling, and comparison with a baseline on the same data for the ACT TT statistic near $1.30$; that number alone does not determine a significance.

\paragraph{IX.4.}
Complete the native recombination physics, reconcile the drag ruler under one background ledger, and use full likelihoods where warranted. A governing equation for the amplitude/tilt/optical-depth triplet remains distinct from a selected formula for each parameter. Helium, two-photon and Lyman-$\alpha$ processes, additional atomic levels, and a nucleosynthesis network require their own physical inputs. Lack of an optimizer does not remove choices already made before a run.

\subsubsection*{Chapter X: many-source organization and halo tests}

\paragraph{X.1--X.3.}
The observed plateau motivated $10/\pi$, so four-decimal agreement is not a prospective prediction. The review's uncertainty estimate uses a large-sample median approximation and the reported scatter; it is illustrative until the actual distribution, weighting, correlation, and systematic model are specified. Use a justified resampling or sampling model for the uncertainty. An untouched sample with frozen selection, mass-to-light assumptions, statistic, and rejection rule can supply the prospective CR034 test. The radial organization law, disk geometry, sensitivity to stellar mass-to-light ratio, profile parameters, and like-for-like radial-acceleration comparison remain necessary.

\paragraph{X.4--X.5.}
A label such as trapped accumulation does not specify a dark component's dynamics. State its compact fraction, mass function, spatial and velocity distributions, baryonic coupling, and formation history. Compact-object, lensing, clustering, and accretion constraints depend on those choices. The heavy-seed ledger and microlensing question remain open; no universal allowed primordial-hole window is assumed here. Many-source nonzero accumulation is not by itself a derived galaxy halo profile.

\subsubsection*{Chapter XI: radiation and ringdown}

\paragraph{XI.1--XI.3.}
The decisive spectral problem needs the actual field operator, causal boundary condition, excitation, and physical units. A reflecting surface, an absorbing horizon, and a dissipative boundary yield different tests, including possible echoes. Binary-pulsar flux requires a moving-source calculation and energy balance; the retained Chapter IV external-contact receipt does not supply that derivation. Spin must be incorporated before observed rotating remnants are used to assess the nonrotating rational candidates. A nonrotating comparator is useful as a controlled theoretical baseline, not a generic observed remnant.

\paragraph{XI.4--XI.6.}
The selection histories of $3/8$ and $4/45$ remain explicit. Broader-grammar counts do not turn either into a mode calculation or a probability of correctness. The $1/8$ radiation ceiling needs a precise initial-energy normalization, permitted source class, and accounting of final mass and all outgoing channels before an event can falsify it. Comparisons with area bounds, numerical-relativity efficiencies, or catalog energies require matched assumptions and uncertainties. A flux law, higher-mode structure, spin dependence, and literature placement remain the necessary continuation.

\subsubsection*{Chapter XII: future road and Home continuation}

\paragraph{XII.1--XII.2.}
The past-light-cone limit $z\to\infty$ does not establish the future-time limit $a\to\infty$. The expression $H_{\infty,\rm readout}=H_0\Omega_\Lambda^{3/2}$ is conditional on the background adapter, proposed future road limit, and definition of the rate. A constant conversion between two distance coordinates does not by itself change their fractional expansion rates; differentiate the actual time-dependent observable before assigning the extra road factor. The model needs an evolution and conservation law, time coordinate, stability analysis, and finite-redshift distance relations. Distance-duality tests require a stated mapping to luminosity and angular-diameter distances and photon transport. Chapter V's source-law conditional must replace an unconditional physical claim of infinite crossing cost.

\paragraph{XII.3--XII.5.}
Home re-indexing remains a conjectural continuation. A transition requires a physical process, conserved quantities, matching conditions for child floor and inventory, and energy/information accounting. A parent-readable consequence must pass through an allowed channel under Chapter VI; hidden child labels alone are not an observable. Primordial-hole formation estimates from horizon mass and temperature depend on expansion history, collapse efficiency, particle degrees of freedom, and the mass distribution. The earlier review's approximate epochs are conditional comparators, not derived SAM formation times or universal constraints. Birkhoff-type statements require their gravitational field-equation hypotheses; they cannot be imported during dynamical collapse solely from the static exterior. Any hard-surface claim must specify its own stress, emission, and causal behavior.

\subsection*{5. A coherent route forward}

The most economical next work reuses complete source calculations rather than replacing them with numerical identities. A source-measure construction could define inventory and road operators on one declared preparation and ask whether both normalize to $A_0R$. A boundary instrument could test the baryon survival factor. A correlated completion calculation could test the proposed sixth power rather than assuming independent writes. These are proposed constructions: neither their correction size nor its sign has been established by this volume.

A physical unit chart can then propagate calibration and experimental uncertainty through existing conditional calculations. It must not disguise model-selection uncertainty as arithmetic precision. A single physically selected boundary law could support both the thermal and ringdown problems, with separate preparation and rate calculations. A time-dependent road law could connect finite-redshift observations to a future limit. It cannot be replaced by repeating the shared number $1/\pi$ in two different roles.

For observational contact, prioritize tests whose outcome was not used to select the formula: a predeclared halo sample, a boundary eigenproblem before mode comparison, light-element abundances from a fixed packet, and complete likelihood comparisons under frozen data and nuisance rules. For gravitational contact, prioritize the spatial metric response, moving-source flux, species universality, and the actual nonuniform source-to-detector cap. These choices do not guarantee a discriminator; they make agreement and failure scientifically legible.

The volume therefore closes with a substantial body of exact and conditional source mathematics, not with a completed derivation of cosmology or microscopic matter. Its strongest continuity is methodological: retain the actual state, all physical complements and payments, the receiving law, and the observable at every handoff. Volume II can supply further matter and calibration candidates, but any dependency running back into Volume I must be recorded explicitly rather than presented as an already independent closure.

\subsection*{Provenance and contribution statement}

Sean Brady is the originator and conceptual director. Chapters I--VI retain their recorded OpenAI ChatGPT and Codex research-collaboration and co-authorship credits, including the particular analytical and verification contributions named in their source registers. The supplied Chapters VII--XII and earlier closing materials retain their Claude/Claude Code contributions where recorded. This replacement synthesis is an OpenAI Codex editorial and scientific-consistency contribution based on the supplied complete working edition of 19 October 2026. It does not reattribute historical derivations, imply that a collaborator executed another's records, or claim a new native computation.

The chapter source registers remain the authorities for exact operands, archived derivations, receipts, and repository versions. External observational numbers requiring renewed checking are identified as such rather than silently promoted to verified results. Historical review classifications and counts are superseded by the scoped register above: source-reported, conditional, and physically open are deliberately kept distinct.
